\documentclass[tikz,border=4pt]{standalone}
\usepackage{ctex}
\usepackage{amsmath}
\usepackage{pgfplots}
\pgfplotsset{compat=1.18}

% p9-02c: 反比例函数单调性示意图——每支上单调（k>0: 第一象限递减，第三象限递减）
\begin{document}
\begin{tikzpicture}
\begin{axis}[
  width=8cm, height=8cm,
  axis lines=center,
  xlabel={$x$}, ylabel={$y$},
  every axis x label/.style={at={(axis cs:4.5,0)}, anchor=west},
  every axis y label/.style={at={(axis cs:0,4.5)}, anchor=south},
  xmin=-4.5, xmax=4.5,
  ymin=-4.5, ymax=4.5,
  xtick={-4,-3,-2,-1,1,2,3,4},
  ytick={-4,-3,-2,-1,1,2,3,4},
  tick style={color=black},
  ticklabel style={font=\small},
  clip=true,
]
  % y = 6/x
  \addplot[blue, thick, domain=1.35:4.4, samples=80] {6/x};
  \addplot[blue, thick, domain=-4.4:-1.35, samples=80] {6/x};

  % 递减箭头（第一象限：从左到右，y减小）
  \draw[->, blue!70, very thick] (axis cs:1.5, 4.0) -- (axis cs:3.0, 2.0) node[right, blue, font=\small] {递减};

  % 递减箭头（第三象限：从左到右，y减小）
  \draw[->, blue!70, very thick] (axis cs:-3.0, -2.0) -- (axis cs:-1.5, -4.0) node[left, blue, font=\small] {递减};

  % 注意：不能跨越x=0连通两支
  \node[right, font=\small] at (axis cs:0.2, 3.0) {$k>0$};
  \node[font=\footnotesize, align=center] at (axis cs:3.2, 4.0) {每支\\单调};
\end{axis}
\end{tikzpicture}
\end{document}
